\documentclass{KPS_assignment}
% ---------For Persian------------
% \newcommand*{\name}{کیوان جمالی}
% \newcommand*{\id}{403209743}
% \newcommand*{\professor}{دکتر شفاهی}
% \newcommand*{\cdate}{2024/12/26}
% \newcommand*{\course}{درس (شماره)}
% \newcommand*{\assignment}{تمرین شماره1}
% ----------For English------------
\newcommand*{\name}{Keivan Jamali}
\newcommand*{\id}{403209743}
\newcommand*{\professor}{Dr. Shafahi}
\newcommand*{\cdate}{1404/01/16}
\newcommand*{\course}{Airport (20\_582)}
\newcommand*{\assignment}{Homework 1}
% ----------START------------------
\begin{document}
\maketitle
\section*{Problem 1}

\subsection*{Problem Description}

Consider a sequence of aircraft that are scheduled to land on a single runway.  
Let $A_i$ denote the \textbf{planned landing time} of aircraft $i$.  

Due to operational uncertainties, the \textbf{actual landing time} is subject to random error:

\[
L_i = A_i + e_i
\]

where the landing error is normally distributed:

\[
e_i \sim \mathcal{N}(0,6)
\]

and errors are assumed to be independent across aircraft.

For safety reasons, the minimum required separation between two consecutive actual landings is:

\[
L_{i+1} - L_i \ge 80
\]

If this condition is violated, the aircraft must perform a go-around, reducing runway capacity.

The objective is to determine the optimal planned spacing between aircraft that maximizes expected runway capacity.

\subsection*{Distribution of Landing Separation}

Define the actual landing time difference:

\[
D_i = L_{i+1} - L_i
\]

Substituting:

\[
D_i = (A_{i+1} - A_i) + (e_{i+1} - e_i)
\]

Assume constant planned spacing:

\[
s = A_{i+1} - A_i
\]

Since $e_i \sim \mathcal{N}(0,6)$ and errors are independent:

\[
e_{i+1} - e_i \sim \mathcal{N}(0,12)
\]

because variances add:

\[
\text{Var}(e_{i+1} - e_i) = 6 + 6 = 12
\]

Therefore:

\[
D_i \sim \mathcal{N}(s,12)
\]

\subsection*{Probability of Safe Landing}

The probability that a landing is successful is:

\[
P(D_i \ge 80)
\]

Using standardization:

\[
P(D_i \ge 80)
=
1 -
\Phi\left(
\frac{80 - s}{\sqrt{12}}
\right)
\]

where $\Phi(\cdot)$ is the standard normal cumulative distribution function.

\subsection*{Capacity Function}

Runway capacity is defined as the expected number of successful landings per unit time.

Since each landing consumes $s$ units of time on average, the expected capacity is:

\[
C(s)
=
\frac{
1 -
\Phi\left(
\frac{80 - s}{\sqrt{12}}
\right)
}{s}
\]

This represents the expected number of successful landings per unit time.

\subsection*{Optimization Problem}

The optimal planned separation $s^*$ is obtained by solving:

\[
\max_{s > 0}
\quad
C(s)
=
\max_{s > 0}
\quad
\frac{
1 -
\Phi\left(
\frac{80 - s}{\sqrt{12}}
\right)
}{s}
\]

Since a closed-form solution does not exist, $s^*$ must be computed numerically.

\subsection*{Final Capacity}

After computing $s^*$ numerically, the maximum runway capacity is:

\[
C^* = C(87.471^*) = 0.011255
\]

If time is measured in seconds, runway capacity in landings per hour is:

\[
\text{Capacity per hour}
=
3600 \times C^* = 40.518
\]

\section*{Problem 2}

Let:

\begin{itemize}
    \item $T \sim \mathcal{N}(\mu, \sigma^2)$ be the total travel + airport processing time.
    \item $H$ be the time between two consecutive flights (in minutes).
    \item $x$ be the decision variable: time before the next flight at which the passenger leaves home.
\end{itemize}

If a flight departs at time $0$, the passenger leaves at time $-x$.
Arrival time at the gate is:

\[
A = -x + T
\]

\subsection*{Waiting Time Function}

Two cases arise:

\[
W(x,T) =
\begin{cases}
x - T & \text{if } T \le x \\
H + x - T & \text{if } T > x
\end{cases}
\]

Equivalently,

\[
W(x,T) = x - T + H \mathbf{1}_{\{T > x\}}
\]

\subsection*{Expected Waiting Time}

Taking expectation:

\[
\mathbb{E}[W(x)]
=
x - \mu
+
H \mathbb{P}(T > x)
\]

Since $T$ is normally distributed,

\[
\mathbb{P}(T > x)
=
1 - \Phi\left(\frac{x-\mu}{\sigma}\right)
\]

Therefore, the objective function is:

\[
f(x)
=
x - \mu
+
H \left[
1 - \Phi\left(\frac{x-\mu}{\sigma}\right)
\right]
\]

\subsection*{Optimization Problem}

\[
\min_{0 \le x \le H}
\;
f(x)
\]

\subsection*{First Order Condition}

Taking derivative:

\[
f'(x)
=
1
-
\frac{H}{\sigma}
\phi\left(\frac{x-\mu}{\sigma}\right)
\]

Setting $f'(x)=0$:

\[
1
=
\frac{H}{\sigma}
\phi\left(\frac{x-\mu}{\sigma}\right)
\]

Thus, the optimal $x^*$ satisfies:

\[
\phi\left(\frac{x^*-\mu}{\sigma}\right)
=
\frac{\sigma}{H}
\]

This implicitly determines the optimal departure time.

\subsection*{Answers to Questions}

\subsubsection*{(a) Optimal $x$ when $H = 3$ hours}

Convert hours to minutes:

\[
H = 3*60 = 180 \text{ minutes}
\]

Solve:

\[
\phi\left(\frac{x^*-72}{\sqrt{17}}\right)
=
\frac{\sqrt{17}}{180}
\]

\[
x_1^* = 65 \text{ minutes}
\]
\[
x_2^* = 79 \text{ minutes}
\]

Since $x$ must be greater than the 72(mean of wait time), the optimal departure time is:

\[
x^* = 79 \text{ minutes}
\]

The mean waiting time is: 15.06 minutes.

\subsubsection*{(b) Effect of Increasing $H$ to 6 Hours}

If $H$ increases:

\[
\phi(z^*) = \frac{\sigma}{H}
\]

Since $H$ increases, the right-hand side decreases.
Because $\phi(z)$ decreases as $|z|$ increases from $0$,

\[
|z^*| \text{ increases}
\]

Therefore,

\[
x^* \text{ increases}
\]

\[
x_1^* = 63.481 \text{ minutes}
\]
\[
x_2^* = 80.518 \text{ minutes}
\]

\textbf{Conclusion:}  
When flights are less frequent, the passenger departs earlier (larger safety buffer) -> 80.518.

The mean waiting time is: 15.51 minutes.

\subsection*{(c) Mean Number of People Waiting in Airport}

Assume:

\begin{itemize}
    \item $D$ destinations
    \item Each destination has flight interval $H$
    \item Demand per destination = $\lambda$ persons per flight
\end{itemize}

Total arrival rate per destination:

\[
\lambda_d = \frac{\lambda}{H}
\]

Using Little's Law:

\[
L = \lambda_d \cdot \mathbb{E}[W(x^*)]
\]

Total number waiting across all destinations:

\[
L_{total} = D \cdot \lambda_d \cdot \mathbb{E}[W(x^*)]
\]

\[
L_{total} = 30 \cdot \frac{60}{180} \cdot \mathbb{E}[W(79)] = 150.60
\]


\section*{Problem 3}

\subsection*{Current situation}
\[
t = \frac{1}{40-q}
\]
t is mean time for service in Operation each hour.

\[
total = \frac{q}{40-q}
\]
total is the service time, in one hour.

\[
day = \frac{24 \times q}{40-q}
\]
all the service time, each day

\[
year = \frac{365 \times 24 \times q}{40-q}
\]
all the service time, each year


\[
cost = \frac{3000 \times 365 \times 24 \times q}{40-q} = 26280000 \times \frac{q}{40 - q}
\]
This is the yearly cost for current situation.


\subsection*{Future situation}
\[
t = \frac{1}{40-(q/2)}
\]
t is mean time for service in Operation each hour.

\[
total = \frac{q}{40-(q/2)}
\]
total is the service time, in one hour for one Band.

\[
day = \frac{24 \times q}{40-(q/2)}
\]
all the service time, each day for one Band.

\[
year = \frac{365 \times 24 \times q}{40-(q/2)}
\]
all the service time, each year for one Band.

\[
year = \frac{2 \times 365 \times 24 \times q}{40-(q/2)}
\]
all the service time, each year for both Band.

\[
cost = \frac{3000 \times 2 \times 365 \times 24 \times q}{40-(q/2)} = 52560000 \times \frac{q}{40 - (q/2)}
\]
This is the yearly cost for future situation.

\[
build = \frac{i \times (i + 1)^{40}}{(i + 1)^{40} - 1} \times 60M + 3M = 8031609.69 (8.0316M)
\]
This is the yearly cost build the new Band.

\[
build = 8031609.69 + 52560000 \times \frac{q}{40 - (q/2)}
\]
The total cost for future in reality.

\subsection*{Compare to find the best year}

\[
q = 20 \times (1 + 0.1) ^ n
\]

Therefore current yearly cost in year n is:
\[
current = 26280000 \times \frac{20 \times (1 + 0.1) ^ n}{40 - (20 \times (1 + 0.1) ^ n)}
\]

Therefore future yearly cost in year n is:
\[
future = 8031609.69 + 52560000 \times \frac{20 \times (1 + 0.1) ^ n}{40 - \frac{(20 \times (1 + 0.1) ^ n)}{2}}
\]

We solve it for n=0:
\[
26.28M \leq 43.07M
\]

We solve it for n=1:
\[
32.12M \leq 47.90M
\]

We solve it for n=2:
\[
40.25M \leq 53.62M
\]

We solve it for n=3:
\[
52.29M \leq 60.45M
\]

We solve it for n=4:
\[
71.80M \leq 68.72M
\]

Therefore, before the start of \textbf{year 4}, we should to build the new Band.

\begin{figure}[h]
    \includegraphics{q3.png}
    \caption{During 7 years, coparison the costs.}
\end{figure}

\end{document}

